\section{Robustness Summary}
\label{sec:robustness}

All findings reported in Sections~\ref{sec:exp1_results}--\ref{sec:exp4b_results} are subjected to a comprehensive suite of robustness checks detailed in Appendix~\ref{sec:appendix_robustness}. This section summarises the key diagnostics.

\subsection{Inference Corrections}

Across all six sub-experiments, the dominant effects survive both Holm--Bonferroni and Benjamini--Hochberg corrections. HC3 heteroskedasticity-consistent standard errors produce negligible changes in significance status, with the fraction of effects changing significance ranging from 0\% (Experiment~4a) to 6.7\% (Experiment~1). Rademacher wild bootstrap $p$-values for the top~10 effects in each response align closely with asymptotic $p$-values, with differences below 0.01 in all cases.

\subsection{Model Adequacy}

LightGBM cross-validated $R^2$ values match or fall below OLS $R^2$ in Experiments~2, 3b, 4a, and~4b, confirming that the linear model with two-way interactions is correctly specified. In Experiments~1 and~3a, LightGBM modestly exceeds OLS for some responses, indicating detectable nonlinear structure; however, the dominant factor rankings (number of bidders, auction type) are unchanged under the nonparametric model. PRESS-based leave-one-out cross-validated $R^2$ gaps remain below 0.10 in all experiments, indicating minimal overfitting.

\subsection{Variable Selection}

LASSO variable selection retains auction type and number of bidders across all six sub-experiments. No interaction survives LASSO whose parent main effects were dropped, confirming effect heredity. LASSO-selected variables are a subset of the OLS-significant effects in every case.

\subsection{Distributional Heterogeneity}

Quantile regression at the 10th, 25th, 50th, 75th, and 90th percentiles confirms that factor effects are broadly stable across the response distribution. The most notable heterogeneity occurs in Experiment~3a, where the auction type coefficient grows from the lower to upper quantiles, indicating that first-price auctions disproportionately suppress revenue among the best-performing configurations. In Experiment~4a, the bidder objective effect grows toward the upper tail, reflecting that the revenue penalty from utility-maximizing objectives is concentrated among high-revenue cells.

\subsection{Supplementary Robustness Checks}

Discretisation sensitivity analysis across grid sizes of 6, 11, and 21 discrete bid levels (Experiments~1--3) confirms that effect rankings are stable across granularities. Budget sensitivity analysis (Experiment~4a) with budget multipliers $m \in \{0.25, 0.5, 1.0\}$ confirms that the number of bidders retains its dominant ranking across all budget levels. Full diagnostic tables, quantile regression figures, and per-experiment model adequacy statistics appear in Appendix~\ref{sec:appendix_robustness}.

All primary findings survive these checks. Section~\ref{sec:discussion} synthesises the validated results across all six sub-experiments.
